% VotaSJ --- Long Technical White Paper
% SBC LaTeX template (sbc-template.cls). See README.md in this folder
% for how to obtain the class file from the SBC website.
\documentclass[12pt]{article}

\usepackage{sbc-template}
\usepackage{graphicx,url}
\usepackage[utf8]{inputenc}
\usepackage[english]{babel}
\usepackage{hyperref}
\usepackage{listings}
\usepackage{xcolor}
\usepackage{booktabs}
\usepackage{array}
\usepackage{microtype}

\hypersetup{
  colorlinks=true,
  linkcolor=black,
  citecolor=black,
  urlcolor=blue!60!black
}

\lstset{
  basicstyle=\ttfamily\small,
  breaklines=true,
  columns=fullflexible,
  frame=single,
  captionpos=b,
  showstringspaces=false
}

\sloppy

\title{VotaSJ: A Blockchain-Anchored Participatory Budgeting Platform for S\~ao Jos\'e/SC}

\author{Davi Ludvig Longen Machado\inst{1} \and
        Lucas Furlanetto Pascoali\inst{1} \\
        Arthur Clasen de Melo\inst{1} \and
        Pedro Henrique Tesman Mansani da Silva\inst{1}}

\address{Departamento de Inform\'atica e Estat\'{\i}stica (INE) \\
  Universidade Federal de Santa Catarina (UFSC) \\
  Florian\'opolis --- SC --- Brasil
  \email{\{davi.ludvig,lucas.furlanetto,arthur.clasen,pedro.tesman\}@grad.ufsc.br}}

\begin{document}
\maketitle

\begin{abstract}
Municipal participatory budgeting (PB) in Brazil suffers from three structural
failures: turnout below 2\,\% of the eligible population, operational cost in
the order of hundreds of thousands of reais per cycle, and a tally that the
City Hall both runs and announces --- so every contested result ends in an
unverifiable accusation of fraud. We propose VotaSJ, a permissioned-write,
permissionless-read application on Polygon PoS that uses gov.br for
eligibility, smart contracts for cycle state and tally, IPFS for proposal
metadata, and an on-chain milestone log for execution accountability.
Administration is held by a 3-of-5 multisig (City Hall \(\times\)3, UFSC,
TCE-SC) and contract upgrades are gated by a 14-day timelock. This paper
argues the narrow technical case for blockchain in this domain, details the
contract design (implementation, 40-test Hardhat suite), reports gas
measurements (\(\approx\)430k gas per full citizen flow), and discusses the
threat model, LGPD posture, and migration path to a Semaphore-based vote-secrecy
expansion. We do not claim that blockchain ``solves democracy''; we claim it
is the right primitive for the narrow set of steps where operator honesty
cannot be assumed.
\end{abstract}

%% Resumo intentionally omitted: VotaSJ project artifacts are English-only
%% per the team's documented language policy.

\section{Introduction and Executive Summary}
\label{sec:intro}

Brazilian municipal participatory budgeting (PB) has been around since the
1990s, but in cities the size of S\~ao Jos\'e/SC it rarely works in
practice. Turnout sits under 2\,\% of eligible voters, a single cycle costs
the city hundreds of thousands of reais in logistics, and the City Hall
both runs the count and announces the winner --- which means every
contested result ends in an accusation of fraud that nobody can prove or
disprove. The institution survives mostly as political theatre.

VotaSJ replaces the parts of that institution where a public ledger
genuinely helps --- the eligibility commitment, the vote, the tally, and
the milestone log of how the winning proposal is executed --- and leaves
everything else off-chain. Citizens authenticate with
\texttt{gov.br}~\cite{govbrmanual}, the federal digital ID; the backend
computes a salted credential hash and commits it to Polygon
PoS~\cite{polygondocs} through a \texttt{VoterRegistry}; the
\texttt{ParticipatoryBudget} cycle contract handles proposals, voting, and
tally; an \texttt{ExecutionTracker} records contract signing, work start,
milestones, and completion of the winning proposal, with IPFS-anchored
evidence. Administration sits on a 3-of-5 Safe (City Hall $\times 3$, UFSC
$\times 1$, TCE-SC $\times 1$), and any contract upgrade is gated by a
14-day OpenZeppelin \texttt{TimelockController}~\cite{openzeppelin5}, so
the rules cannot change inside an active cycle.

Why blockchain at all? Because no single party in this system is trusted
by all the others. Citizens distrust the City Hall to count its own vote;
the opposition distrusts the incumbent; auditors distrust the speed of
paper reconstruction. What the system has to produce is a count that
nobody has to take on faith, and a public L2 with an external validator
set produces that property without anyone in S\~ao Jos\'e having to be
honest. Sections~\ref{sec:problem}--\ref{sec:why-bc} make that case;
\ref{sec:arch}--\ref{sec:dataflow} describe the system;
\ref{sec:security}--\ref{sec:scale} cover security, privacy and scale;
\ref{sec:workflow}--\ref{sec:economics} report what it costs;
\ref{sec:team}--\ref{sec:gtm} introduce the team and go-to-market plan;
\ref{sec:roadmap}--\ref{sec:conclusion} close with roadmap, risks, and
conclusion.

What VotaSJ is \textbf{not}: an attempt to ``tokenize democracy'' or to
replace deliberation with code. Citizens pay nothing. There is no fungible
vote-token. Deliberation, mobilization, and assembly happen off-chain
where they belong. The blockchain is just the record of who is allowed to
vote, who voted for what, and what happened to the money --- the four
things the institution currently cannot prove.

\section{Problem Statement}
\label{sec:problem}

S\~ao Jos\'e/SC has approximately \textbf{180{,}000 eligible voters} and runs
(irregularly) a participatory-budget process inherited from the Porto Alegre
school of the early 1990s, framed legally by the \emph{Estatuto da
Cidade}~\cite{estatutodacidade}. Field data from comparable mid-size Brazilian
municipalities indicates:

\begin{itemize}
  \item \textbf{Turnout below 2\,\%} in the last three documented cycles ---
        under 3{,}600 actual participants decide the allocation of tens of
        millions of reais of investment-margin budget.
  \item \textbf{Operational cost per cycle of BRL 300k--500k}, dominated by
        physical assembly logistics, poll-worker payments, manual tallying,
        and administrative reconciliation.
  \item \textbf{A tally that takes \(\approx\)3 weeks} of paper reconstruction
        after the last assembly closes, during which the legitimacy of the
        result is essentially indefensible.
  \item \textbf{No execution feedback loop}: even when a proposal wins,
        citizens have no programmatic way to verify whether the committed
        budget was actually executed against the winning project.
\end{itemize}

Three additional second-order effects compound the failure: (i)
\textit{self-selection bias} because evening assemblies systematically
exclude shift workers, parents of small children, and the geographically
peripheral population; (ii) \textit{fraud asymmetry}, since the City Hall
operates, counts, and announces the result; and (iii) \textit{bureaucratic
opacity post-vote}, since the winning proposal disappears into the municipal
financial system (SIG) after the cycle closes.

\subsection{Why existing digital systems fail}
\label{sec:digital-fail}

Since roughly 2015 Brazilian municipalities have experimented with digital PB
tools --- from off-the-shelf survey platforms (Decidim, Consul, SurveyMonkey)
to in-house web portals integrated with municipal SSO. These deployments
share four failure modes:

\begin{itemize}
  \item \textbf{Single-operator trust.} The City Hall owns the database, the
        application server, and the announcement. A determined operator can
        modify results without detectable trace.
  \item \textbf{No public auditability.} Citizens get an HTML page; auditors
        get a CSV upon request weeks later. Neither version is independently
        verifiable against the original event stream because the event stream
        is, in principle, a SQL table the operator can rewrite.
  \item \textbf{No Sybil resistance.} Most deployments fall back to CPF
        \(+\) email confirmation, which is not Sybil-resistant.
  \item \textbf{No execution accountability.} None of the deployed systems
        link the vote result back to budget execution.
\end{itemize}

VotaSJ does not claim that blockchain solves all four problems by default.
It claims, more narrowly, that \textbf{blockchain solves the first two
structurally}; that the third is solvable by anchoring identity to
\texttt{gov.br}, an existing federal credential infrastructure the operator
does not control; and that the fourth is solvable by reusing the same
substrate that records the vote to record the milestone.

\section{Why Blockchain (and Why That Is Not a Tautology)}
\label{sec:why-bc}

Before settling on a blockchain we considered three simpler alternatives:
a normal web application backed by a relational database, the same setup
with a cryptographic audit log on top (Merkle roots, signed receipts), and
a PKI-only design where each voter holds a certificate and votes are
detached signatures. Each one fixes part of the problem and leaves another
part open.

\subsection{Traditional web servers and relational databases}

This is what every existing Brazilian deployment uses. Integrity depends
entirely on operator honesty; the auditor has access only to what the
operator chooses to expose. Adequate for \textit{non-political} civic
systems; unfit for one whose entire value proposition is convincing a
skeptical opposition that the count is honest.

\subsection{Centralized DB with cryptographic audit log}

The operator publishes Merkle roots of the event stream and citizens receive
signed receipts. A meaningful improvement over §3.1, and the strongest
non-blockchain candidate (the architecture behind Certificate Transparency).
The auditor can prove \textit{inclusion} and \textit{exclusion} of a given
receipt against the published root, but \textbf{cannot prove completeness}
--- there may be additional events the operator chose not to commit. The
``completeness'' property is exactly the one a political adversary will
attack.

\subsection{PKI-only systems}

Every voter holds a certificate; votes are detached signatures. Per-vote
integrity is excellent, but \textit{aggregation integrity} is nil --- the
aggregator still decides what counts as the final tally. Operational reality
also dooms it: governmental PKI rollout (e.g., ICP-Brasil) has been
notoriously expensive and slow.

\subsection{Public blockchain}

Integrity is enforced by a consensus protocol on a validator set the
operator does not control~\cite{wood2014yellowpaper}; auditability is
trivial; completeness is by construction; live verification is continuous
with \(\approx\)2\,s finality on Polygon PoS. A blockchain is the only
candidate in this list where the integrity guarantee does not collapse if
the operator becomes adversarial.

\subsection{On-chain only where it matters}
\label{sec:on-chain-only}

\noindent\begin{tabular}{p{0.45\textwidth} c p{0.32\textwidth}}
\toprule
\textbf{Subproblem} & \textbf{On-chain?} & \textbf{Why} \\
\midrule
Proposal metadata (title, images) & no & IPFS; commit only the CID. \\
Personal data (CPF, name, address) & no & LGPD~\cite{lgpd}; never on-chain. \\
Proving eligibility & partial & \texttt{gov.br} proof, hash on-chain. \\
Casting a vote & yes & Operator honesty cannot be assumed. \\
Counting votes & yes & A \texttt{closeCycle()} on-chain is the only count nobody must trust. \\
Recording the winner & yes & Same. \\
Tracking budget execution & yes & The accountability loop the institution fails to close. \\
Notifications, PDFs, charts & no & Standard off-chain web. \\
\bottomrule
\end{tabular}

\section{System Architecture}
\label{sec:arch}

VotaSJ is a three-layer system. Only the on-chain layer is authoritative;
the off-chain backbone is a best-effort cache and integration; the citizen
layer is the UX and the audit surface.

\begin{lstlisting}[caption={High-level architecture}]
+--------------------------------------------------------------+
| L1 - Citizen layer                                           |
|   * Progressive Web App (React + ethers v6)                  |
|   * Assisted kiosks in CRAS / UPAs / public libraries        |
|   * Public audit portal (read-only subgraph queries)         |
+--------------------------------------------------------------+
| L2 - Off-chain backbone                                      |
|   * Voter Registry DB (PostgreSQL, LGPD-controlled)          |
|   * gov.br OIDC relay (Serpro federated identity)            |
|   * IPFS pinning (proposal metadata)                         |
|   * The Graph subgraph (public indexer)                      |
|   * Chainlink Automation (cycle-closing keeper)              |
|   * Notification worker (push/email/SMS - non-authoritative) |
+--------------------------------------------------------------+
| L3 - On-chain substrate of record (Polygon PoS)              |
|   * VoterRegistry.sol                                        |
|   * ParticipatoryBudget.sol (cycle, propose, vote, close)    |
|   * ExecutionTracker.sol (milestone commitments)             |
|   * TimelockController + 3-of-5 Safe multisig                |
|   * UUPS proxy pattern (upgradeable behind timelock)         |
+--------------------------------------------------------------+
\end{lstlisting}

\subsection{Why Polygon PoS, not Ethereum L1, not a private chain}

Ethereum L1 is the natural conservative choice but at current gas prices it
produces a per-vote cost in the order of USD 0.20--1.00, which over 50{,}000
votes per cycle becomes prohibitive without buying additional security
relative to a major L2. A private/permissioned chain (e.g., Hyperledger
Fabric) fails the central test: the operator (City Hall) would also control
the validator set --- the chain degenerates into the configuration ruled out
in §3.2. A roll-your-own validator set (UFSC + TCE-SC + Public Prosecutor's
Office running their own nodes) is plausible but requires significant
institutional negotiation and is fragile against political reconfiguration;
it is kept as a deliberate fallback. Polygon PoS gives us a public L2
with \(\approx\)100 external validators, \(\approx\)2\,s finality, EVM
equivalence, and per-tx fees in fractions of a cent~\cite{polygondocs}.

\section{Smart Contract Design}
\label{sec:contracts}

VotaSJ splits the on-chain state into two contracts with deliberately
narrow responsibilities --- a \texttt{VoterRegistry} and a
\texttt{ParticipatoryBudget} cycle machine --- plus an
\texttt{ExecutionTracker} for milestone commitments. The contract suite is
exercised by a 40-test Hardhat suite~\cite{hardhatdocs}.

\subsection{VoterRegistry}

Maintains the set of eligible voter addresses keyed by a salted credential
hash. Admin-only writes (\texttt{registerVoter}, \texttt{revokeVoter}),
public reads. The registry is small (one address per citizen) and rarely
changes.

\subsection{ParticipatoryBudget}

A finite-state cycle machine with three states (Pending, Open, Closed).
Public writes are gated on \texttt{registry.isRegistered(msg.sender)};
admin-only writes (\texttt{openCycle}) are gated on the multisig. The
implementation is wrapped behind a UUPS proxy~\cite{openzeppelin5}.

State transitions:

\begin{lstlisting}
Pending --openCycle()--> Open
Open    --closeCycle() after closesAt--> Closed
Closed  --openCycle()--> Open  (next cycle)
\end{lstlisting}

Three invariants are enforced at the contract level and exercised in the
test suite:

\begin{enumerate}
  \item Only registered voters can submit a proposal or cast a vote.
  \item One vote per voter per cycle.
  \item No cycle reopens once closed.
\end{enumerate}

Tally is deterministic: \texttt{closeCycle()} iterates the proposals of the
current cycle and emits \texttt{CycleClosed(cycleId, winningProposalId,
winningVotes)}. There is no off-chain tally service to trust.

\subsection{ExecutionTracker}

Once a cycle closes with a winning proposal, the City Hall financial system
(SIG) consumes the on-chain event via a webhook and triggers the off-chain
execution flow. As budget is committed and disbursed, milestone commitments
are written back via \texttt{ExecutionTracker.recordMilestone(proposalId,
milestoneId, cidEvidence)} by a multisig-controlled relayer, with
IPFS-anchored evidence (invoices, photographs, contract scans). This is the
\textit{accountability loop} that the existing institution fails to close.

\subsection{Storage layout and gas economy}

\noindent The packing decisions are deliberate: \texttt{opensAt} and
\texttt{closesAt} are \texttt{uint64} (sufficient until year 292277026596),
letting the EVM pack them into a single storage slot, while larger fields
land in their own slots. Empirical gas measurements from the Hardhat test
suite indicate roughly \(\approx\)430{,}000 gas for a full citizen flow
(register \(+\) submit \(+\) vote \(+\) read). At a
Polygon gas price of 30\,gwei and a MATIC price of \(\approx\)USD 0.40,
the per-citizen-per-cycle cost lands at \(\approx\)USD 0.005 --- three
orders of magnitude below the existing process cost.

\subsection{Events as the indexing surface}

We design events as a first-class API, not a debugging afterthought:
\texttt{CycleOpened}, \texttt{ProposalSubmitted}, \texttt{VoteCast},
\texttt{CycleClosed}. The public subgraph~\cite{thegraphdocs} is built
strictly from these events; the audit portal queries the subgraph. Citizens,
journalists, TCE-SC, and the Public Prosecutor's Office consume the same
surface --- no one gets a privileged audit channel because no one needs one.

\subsection{Upgradeability (UUPS + Timelock)}

The production deployment wraps \texttt{ParticipatoryBudget} behind a UUPS
proxy. The implementation contract is the only target of \texttt{upgradeTo},
which is restricted to the multisig and gated by a 14-day
\texttt{TimelockController}~\cite{openzeppelin5}. UUPS over Transparent
Proxy reduces per-call gas overhead by \(\approx\)2{,}000\,gas; the 14-day
timelock is calibrated to be longer than the critical voting interval. A
second guardrail is enforced at the application level: \texttt{openCycle()}
reverts if a pending upgrade is queued. The combination makes
``rewrite-the-tally-rules-during-the-vote'' a strict impossibility under
the multisig threat model.

\subsection{Implementation notes for the next developer}

This subsection collects the design decisions that are not obvious from
the contract source. We document them here because each is a place where
a future maintainer is likely to second-guess us.

\begin{itemize}
  \item \textbf{Revert error strings follow a \texttt{"Contract: reason"}
        convention} (\texttt{"Budget: not admin"},
        \texttt{"Budget: cycle not open"}, \texttt{"Registry: zero
        address"}). The frontend pattern-matches on the prefix to route
        errors to the right UI surface; renaming a contract changes the
        prefix and is a backward-incompatible UI change.
  \item \textbf{\texttt{opensAt} and \texttt{closesAt} are
        \texttt{uint64}} (sufficient through year 292277026596) so they
        pack into a single storage slot with \texttt{status} and
        \texttt{proposalCount}. Using \texttt{uint256} everywhere would
        cost an extra slot per cycle and \(\approx\)3{,}000\,gas per
        \texttt{openCycle} for no practical benefit.
  \item \textbf{The \texttt{VoterRegistry} reference is
        \texttt{immutable}} in \texttt{ParticipatoryBudget}. The registry
        address is fixed at construction; rebinding to a different
        registry would require redeploying the budget contract and
        migrating cycle history. This is deliberate: a swappable registry
        is a larger attack surface than a redeployment is a maintenance
        cost.
  \item \textbf{Single active cycle invariant.} \texttt{openCycle}
        requires the previous cycle to be \texttt{Closed}. Multi-cycle
        parallelism (one per administrative region) is a roadmap item.
  \item \textbf{\texttt{closeCycle} iterates linearly} over
        \texttt{proposalCount} to find the winner. We bound
        \texttt{proposalCount} at 500 per cycle; above that the close
        transaction approaches Polygon's 30M block gas limit. For larger
        deployments, the tally moves to an off-chain ZK rollup with a
        single on-chain verification --- that work is scoped but not yet
        built.
  \item \textbf{Vote receipts are events, not return values.} Callers
        reconstruct receipts from the \texttt{VoteCast} event, which lets
        the subgraph be the single source of truth for ``did my vote land''.
  \item \textbf{No comments, ratings, or rich proposal interaction on-chain.}
        Anything that could carry personal data, opinion text, or
        moderation surface stays in IPFS or in the off-chain backend; the
        contract knows only proposal IDs, vote counts, and milestone hashes.
\end{itemize}

\section{Identity Verification Flow}
\label{sec:identity}

VotaSJ does not authenticate citizens directly. It delegates to
\texttt{gov.br}~\cite{govbrmanual}, the Brazilian federal OIDC identity
provider operated by Serpro:

\begin{enumerate}
  \item The citizen opens the VotaSJ PWA and clicks ``Register to vote''.
  \item The PWA redirects to gov.br with
        \texttt{scope=openid profile govbr\_confianca\_alto}.
  \item gov.br returns a signed assertion containing the citizen's CPF and
        confidence level.
  \item The backend computes \texttt{credentialHash =
        keccak256(cpf || municipality\_code || cycle\_salt || pepper)}, where
        \texttt{pepper} is a server-side secret rotated every cycle. The
        CPF never leaves the boundary; only the hash does.
  \item The backend asks the citizen to bind a wallet address (in-app
        ERC-4337~\cite{eip4337} smart-account, or external).
  \item The wallet \(\leftrightarrow\) \texttt{credentialHash} binding is
        committed via \texttt{VoterRegistry.registerVoter} from the multisig.
  \item From this point, the citizen votes from the wallet. The chain knows
        nothing about the CPF.
\end{enumerate}

The pepper-and-salt construction means the same CPF produces a different
\texttt{credentialHash} each cycle, which prevents long-term linkability of
a voter across cycles by anyone who later obtains the database. The pepper
is held by the City Hall under multisig escrow; it is never required to be
known publicly to verify the integrity of the count.

\textbf{Trust assumption.} We assume gov.br is honest about the identities
it asserts. This is the same assumption every federal service makes; we do
not improve it, but we also do not weaken it. Mitigation against a gov.br
compromise is institutional, not cryptographic --- a credential-hash anomaly
(e.g., a sudden burst of registrations from a single IP block) is detectable
by the auditor multisig and triggers a pause.

\section{Multisig Governance}
\label{sec:governance}

Administrative authority is held by a \textbf{3-of-5 Gnosis Safe}:

\begin{center}
\begin{tabular}{l l}
\toprule
\textbf{Signer} & \textbf{Role} \\
\midrule
City Hall \#1 (Planning Secretariat) & Executive \\
City Hall \#2 (Ombudsman)            & Internal control \\
City Hall \#3 (Legal Department)     & Procedural validity \\
UFSC INE / Blockchain Lab           & Academic auditor \\
TCE-SC (State Court of Audit)        & Institutional auditor \\
\bottomrule
\end{tabular}
\end{center}

The choice is deliberate. A pure City Hall multisig would not solve the
political-asymmetry problem of §\ref{sec:problem}. A purely institutional
multisig would make routine operations slow and politically fraught. The
3-of-5 with three internal and two external signers, threshold 3, encodes:

\begin{itemize}
  \item The City Hall \textbf{cannot} unilaterally open, close, or modify a
        cycle (needs at least one external signer).
  \item The City Hall \textbf{can} still run routine operations (3 internal
        signatures) when no auditor objects.
  \item Any single external signer can veto by withholding their signature.
\end{itemize}

The Safe is itself wrapped by an OpenZeppelin
\texttt{TimelockController}~\cite{openzeppelin5} for any operation that
changes contract code, role assignments, or the signer set. Routine cycle
operations are not timelocked --- they happen at the speed of multisig
coordination.

\section{Data Flow and Processing}
\label{sec:dataflow}

End-to-end, the lifecycle of a single vote:

\begin{enumerate}
  \item Citizen authenticates with gov.br~\cite{govbrmanual}.
  \item Backend issues \texttt{credentialHash}; multisig calls
        \texttt{VoterRegistry.registerVoter}.
  \item Citizen browses proposals; metadata loaded from IPFS, addressed by
        the on-chain CID.
  \item Citizen calls \texttt{ParticipatoryBudget.vote(proposalId)} via an
        ERC-4337~\cite{eip4337} paymaster funded by the City Hall --- the
        citizen pays no gas.
  \item Polygon validators include the transaction; finality in \(\approx\)2\,s.
  \item The subgraph indexer~\cite{thegraphdocs} ingests the
        \texttt{VoteCast} event.
  \item The public audit portal updates the per-proposal count in real time.
  \item After \texttt{closesAt}, Chainlink
        Automation~\cite{chainlinkautomation} calls \texttt{closeCycle()};
        tally is deterministic.
  \item The winning proposal is read by the City Hall's SIG via webhook.
  \item As execution progresses, milestones are committed to
        \texttt{ExecutionTracker}.
\end{enumerate}

The citizen sees a normal mobile UX; the auditor sees a public event log;
the City Hall sees a SIG integration. Nobody is forced to interact with the
blockchain directly except as the system-of-record.

\section{Security Model and Threat Analysis}
\label{sec:security}

\subsection{Threat actors and assets}

\begin{center}
\begin{tabular}{l p{0.32\textwidth} p{0.28\textwidth}}
\toprule
\textbf{Actor} & \textbf{Capability} & \textbf{Asset at risk} \\
\midrule
Adversarial City Hall & 3/5 signers, runs backend & Tally outcome \\
External attacker & Internet, backend compromise & Voter Registry, gas budget \\
Vote buyer & Money, social pressure & Vote secrecy \\
Sybil attacker & Multiple CPFs & Vote weight \\
Polygon validator collusion & \(\geq\)2/3 validators & Liveness, finality \\
Smart-contract bug & --- & Cycle integrity \\
\bottomrule
\end{tabular}
\end{center}

\subsection{Sybil resistance}

The chain itself does not provide Sybil resistance. We obtain it externally
via gov.br: each CPF produces only one \texttt{credentialHash} per cycle, and
each \texttt{credentialHash} is bound to at most one address by the registry.
The remaining vector --- one human, multiple CPFs --- requires either a
stolen identity (separately a crime, detectable by gov.br's anti-fraud
telemetry) or a corrupted gov.br assertion (federal-level compromise, out of
scope). We explicitly do \textbf{not} rely on captcha, phone-number
verification, or email confirmation for Sybil resistance --- none of those
are Sybil-resistant at scale.

\subsection{Admin abuse resistance}

Three layers: (1) critical state changes require 3-of-5 signatures including
\(\geq\)1 external signer; (2) critical \emph{rule} changes additionally
require a 14-day timelock; (3) \texttt{openCycle} reverts if any contract
upgrade is pending --- closing the loophole of ``upgrade now, execute against
the new rules immediately''.

\subsection{Censorship resistance}

The contract is deployed on a public L2 with \(\approx\)100 external
validators. The City Hall cannot unilaterally take VotaSJ offline; doing so
would require either shutting down Polygon (which they do not control) or
pushing a contract upgrade to disable voting (timelocked, multisig-gated,
publicly visible).

\subsection{Smart-contract attack surface}

We treat the smart contract as the single highest-value target:

\begin{itemize}
  \item External audit (OpenZeppelin or Hacken) before every major release,
        budgeted at BRL 40k--80k per audit.
  \item Reentrancy protection via \texttt{ReentrancyGuard}~\cite{openzeppelin5}
        on any state-mutating function that touches an external contract.
  \item Integer overflow safety automatic in Solidity \(\geq\)0.8.
  \item Pausable escape hatch wired to the multisig.
  \item No \texttt{delegatecall} outside the UUPS proxy slot.
  \item Slither + Mythril static analysis in CI on every PR.
\end{itemize}

\subsection{Emergency procedure}

If the multisig detects an active attack (Sybil wave, smart-contract bug,
Polygon-level reorg): (1) any 3-of-5 signers call \texttt{pause()}; (2) the
Public Prosecutor's Office and the citizenry are notified via the public
portal banner; (3) diagnosis is off-chain; the cycle is closed and re-opened
with new parameters via multisig + timelock, or allowed to expire and its
result discarded with public, signed justification.

\section{Privacy and LGPD Compliance}
\label{sec:privacy}

LGPD~\cite{lgpd} is binding on the City Hall as data controller. VotaSJ's
posture:

\begin{itemize}
  \item \textbf{No personal data on-chain.} Ever. CPF, full name, address,
        phone, email --- all live in the off-chain Voter Registry DB.
  \item \textbf{Only credential hashes on-chain.} One-way function;
        unrecoverable without the per-cycle pepper.
  \item \textbf{Per-cycle pepper rotation.} Long-term linkability across
        cycles is prevented at the chain level.
  \item \textbf{Right-to-be-forgotten (LGPD Art. 18).} The on-chain
        \texttt{credentialHash} cannot be deleted but, because it is a
        one-way function of a CPF the system no longer holds, the on-chain
        residue is, in our reading, non-identifying. This interpretation
        has not yet been confirmed by the ANPD for our specific deployment.
  \item \textbf{DPO and Data Processing Agreement.} The City Hall's appointed
        DPO co-signs the LGPD impact assessment.
\end{itemize}

\subsection{Future zk-proof eligibility (Semaphore expansion)}
\label{sec:zk}

The current registry-based design exposes one cryptographic linkage that
the next iteration breaks: an external observer who controls the Voter
Registry DB and observes a \texttt{VoteCast} event can correlate the wallet
address with the citizen. This is not a leak to the \emph{public} (the
wallet is pseudonymous) but it is one to the \emph{City Hall}. The next
iteration replaces the direct \texttt{registry.isRegistered(msg.sender)}
check with a Semaphore-style group-membership
zk-proof~\cite{semaphoredocs,semaphoregithub}:

\begin{itemize}
  \item The eligible cohort is committed as a Merkle root of credential
        hashes (one root per cycle).
  \item The citizen generates a zk-SNARK proving ``I know a leaf in this
        tree and I have not voted before'' without revealing which leaf.
  \item A nullifier deterministic in \((\texttt{credentialHash},
        \texttt{cycleId})\) is published on-chain to enforce
        one-vote-per-citizen without linkability to the registry.
\end{itemize}

Benchmarks of comparable Semaphore-on-Polygon deployments suggest a per-vote
zk-proof generation of \(\approx\)3--5\,s on a mid-range mobile device and an
on-chain verification cost of \(\approx\)250k\,gas, still within the per-vote
economic envelope. The zk migration is scheduled for the second annual
cycle.

\section{Scalability, Performance, and Infrastructure}
\label{sec:scale}

\subsection{Throughput envelope}

Polygon PoS sustains \(\approx\)65\,TPS of complex EVM transactions and
\(\approx\)1--2\,kTPS of simple ones, with finality in the \(\approx\)2\,s
range~\cite{polygondocs}. A peak voting day under the VotaSJ cycle profile
(\(\approx\)20{,}000 citizens voting in the last 48\,h of a 14-day cycle)
implies a peak sustained load of \(\approx\)0.1\,TPS --- three orders of
magnitude below the chain's capacity.

\subsection{Gas-cost considerations}

\noindent\begin{tabular}{l r}
\toprule
\textbf{Operation} & \textbf{Gas (approx.)} \\
\midrule
\texttt{registerVoter} & 80k \\
\texttt{submitProposal} & 120k \\
\texttt{vote} & 95k \\
\texttt{closeCycle} (per proposal) & 30k \\
\bottomrule
\end{tabular}

At 30\,gwei and MATIC \(\approx\) USD 0.40, the per-vote cost is on the
order of USD 0.001. Even with a 10\(\times\) safety margin and ERC-4337
paymaster overhead~\cite{eip4337}, the per-cycle gas budget at 50{,}000
voters fits comfortably under BRL 5{,}000.

\subsection{closeCycle loop bound}

The current \texttt{closeCycle} implementation iterates over all proposals.
We bound the maximum number of proposals per cycle to 500. At
\(500 \times 30\)k\,gas, the close transaction lands at \(\approx\)15M\,gas
--- well under Polygon's 30M block gas limit. For larger municipalities, the
next iteration will move the tally to an off-chain ZK-rollup with on-chain
verification.

\subsection{Infrastructure requirements}

\begin{itemize}
  \item Dedicated Polygon node (Erigon, 8 vCPU / 32\,GB RAM / 4\,TB SSD), with
        a redundant hot standby.
  \item IPFS pinning via a managed provider (Pinata or local Kubo cluster)
        with 10\,GB of pinned content per cycle and 99.9\,\% availability SLA.
  \item The Graph subgraph on the decentralized
        network~\cite{thegraphdocs} with a secondary self-hosted indexer.
  \item Chainlink Automation subscription, funded once per
        cycle~\cite{chainlinkautomation}.
\end{itemize}

\subsection{Reliability and fault tolerance}

Three independent failure scopes:

\noindent\begin{tabular}{l p{0.27\textwidth} p{0.38\textwidth}}
\toprule
\textbf{Scope} & \textbf{Failure} & \textbf{Mitigation} \\
\midrule
Polygon & Reorg, multi-block halt & Multisig pause; cycle replay; ultimate fallback to Ethereum L1 redeployment (contracts are EVM-portable). \\
Backend & gov.br outage, IPFS unpin, indexer lag & Retries; alternative IPFS gateway; subgraph re-sync. \\
Operator & City Hall non-cooperation & The chain keeps running; 2 external signers; any RPC works. \\
\bottomrule
\end{tabular}

The trust-critical path (cast vote, count vote, publish result) only depends
on Polygon being live. Every other failure mode degrades UX, not integrity.

\section{Operational Workflow}
\label{sec:workflow}

A canonical cycle, end-to-end:

\begin{lstlisting}
T-90  City Hall + UFSC + TCE-SC agree on cycle parameters.
T-60  Multisig queues openCycle() in TimelockController (14-day delay).
T-46  openCycle() executes. Cycle is Pending.
T-30  Proposal-submission window opens.
T-16  Proposal window closes; voting window opens.
T-2   Last 48h: peak voting load.
T0    closesAt reached. Chainlink Automation calls closeCycle().
T+0   CycleClosed event emitted; portal shows winner.
T+1   City Hall SIG receives webhook; budget execution kicks off.
T+30/60/90  ExecutionTracker.recordMilestone() with IPFS evidence.
T+365 Next cycle: GOTO T-90.
\end{lstlisting}

No human is required to perform a count, publish a result, or open a
notification thread; all are produced as side-effects of on-chain events.

\subsection{Operational runbook (who does what during a cycle)}

The contract suite removes humans from the count, but the cycle still
requires named operators with explicit handoffs. We document them here so
the team running the first municipal pilot has a checklist rather than a
vibe.

\begin{center}\small
\begin{tabular}{l l p{0.32\textwidth}}
\toprule
\textbf{Stage} & \textbf{Owner} & \textbf{Activity} \\
\midrule
Cycle config & Planning Secretariat & Drafts cycle params; circulates to UFSC + TCE-SC for review. \\
Timelock queue & Secretariat + \(\geq\)1 external signer & Queues \texttt{openCycle} in the \texttt{TimelockController}. \\
Mobilization & Comms + neighborhood associations & PWA, press, kiosks; trains kiosk attendants. \\
Identity onboarding & Backend on-call & Operates the gov.br relay; monitors registration anomalies. \\
Voting window & Backend on-call + multisig & Monitors gas, paymaster, Polygon block production. \\
Close & Chainlink Automation & Calls \texttt{closeCycle()} at \texttt{closesAt}. \\
Result publication & Frontend + comms & Portal updates from subgraph; press release triggered by \texttt{CycleClosed}. \\
Execution tracking & Secretariat + multisig & \texttt{ExecutionTracker.recordMilestone()} with IPFS evidence at each stage. \\
\bottomrule
\end{tabular}
\end{center}

The single biggest operational risk is \textbf{paymaster underfunding}
during the last 48 hours of voting: at 30\,gwei and the gas profile of
\S\ref{sec:scale}, the paymaster needs roughly BRL 2{,}000--3{,}000 of
MATIC held in advance to cover 50k votes with a 5\(\times\) safety margin.
We pre-fund the paymaster two weeks before the voting window and check
the balance daily during the cycle.

The second biggest risk is a \textbf{multisig signer being unreachable}
in an emergency. The 3-of-5 threshold gives one signer of slack, but only
one. The signer set has a documented escalation chain kept off-chain in
the operations handbook.

\section{Economic Feasibility}
\label{sec:economics}

\subsection{Cost reduction vs. traditional process}

\noindent\begin{tabular}{l r r}
\toprule
\textbf{Item} & \textbf{Traditional} & \textbf{VotaSJ} \\
\midrule
Assembly logistics & 60--100k & 0 \\
Poll workers, transport & 80--150k & 5--10k \\
Manual tallying & 40--80k & 0 \\
Administrative reconciliation & 30--50k & 0 \\
Communications, mobilization & 50--80k & 40--60k \\
External smart-contract audit (amortized) & --- & 40--80k \\
Gas + infrastructure & --- & 5--10k \\
\midrule
\textbf{Total (BRL/cycle)} & \textbf{260--460k} & \textbf{90--160k} \\
\bottomrule
\end{tabular}

Net savings: BRL 170k--300k per cycle, conservatively, before counting the
political value of an unimpeachable result.

\subsection{Revenue model}

VotaSJ is a B2G SaaS sold to the City Hall: annual subscription BRL
180--250k, per-cycle setup fee BRL 40k, implementation consulting for
additional municipalities BRL 60--120k. Citizens never pay.

\section{Team}
\label{sec:team}

VotaSJ was designed and implemented by four senior students of the
\textbf{Departamento de Inform\'atica e Estat\'{\i}stica (INE)} at UFSC,
enrolled in \textbf{INE5458 --- Blockchain and Cryptocurrency Technologies}
(2026/1). Each author contributed across the full stack; the primary
specializations are summarized below.

\begin{center}
\begin{tabular}{l p{0.55\textwidth}}
\toprule
\textbf{Member} & \textbf{Primary contribution} \\
\midrule
Davi Ludvig Longen Machado & Smart-contract architecture: cycle state machine, UUPS-proxy and multisig wiring, OpenZeppelin integration, Hardhat test harness. \\
Lucas Furlanetto Pascoali & Smart-contract security: threat modeling, Slither/Mythril CI gates, gas profiling, reentrancy and access-control review. \\
Arthur Clasen de Melo & Off-chain backbone: gov.br OIDC relay design, IPFS pinning topology, The Graph subgraph schema, Chainlink Automation upkeep. \\
Pedro H.\ T.\ Mansani da Silva & Citizen-facing layer: PWA (React + ethers v6), ERC-4337-based wallet abstraction~\cite{eip4337}, accessibility and kiosk flows. \\
\bottomrule
\end{tabular}
\end{center}

The team writes Solidity, TypeScript, and Python daily; all four members
have prior Brazilian-civic-tech exposure through coursework, hackathons,
and contributions to UFSC research projects. The work is positioned within
the long-running research line of the \textbf{UFSC Blockchain Lab} at INE,
which provides academic supervision and acts as the proposed external
auditor signer of the production multisig. External smart-contract
auditing for the first municipal pilot is contracted to a specialized firm
(OpenZeppelin or Hacken) before any binding deployment --- the team
explicitly does not self-audit security-critical code.

\section{Go-to-Market and Marketing Strategy}
\label{sec:gtm}

VotaSJ is a B2G product. Its commercial cycle is slow, reference-driven,
and institutionally gated --- not a consumer-acquisition campaign. The
marketing plan reflects that reality and is built on top of the
\emph{Channels} and \emph{Customer Relationships} blocks of the Business
Model Canvas.

\subsection{Sales motion (long-cycle B2G)}

\begin{itemize}
  \item \textbf{Anchor sale.} A municipal contract with S\~ao Jos\'e/SC
        as the reference deployment. In Brazilian municipal procurement,
        the first signature is the gating asset.
  \item \textbf{Reference-based expansion.} From S\~ao Jos\'e, the next
        addressable cluster is Greater Florian\'opolis (Palho\c{c}a,
        Bigua\c{c}u, Florian\'opolis), followed by mid-size Santa Catarina
        municipalities via \textbf{FECAM} (Federa\c{c}\~ao Catarinense de
        Munic\'{\i}pios) and \textbf{ABM} (Associa\c{c}\~ao Brasileira de
        Munic\'{\i}pios) channels.
  \item \textbf{Institutional dual-use.} TCE-SC and the Public Prosecutor's
        Office function simultaneously as auditor multisig signers \emph{and}
        as the strongest endorsement signal a municipal buyer can receive.
        Their participation in governance is itself the marketing message.
\end{itemize}

\subsection{Communication channels (mapped from the BMC)}

\begin{center}\small
\begin{tabular}{l l l}
\toprule
\textbf{Channel} & \textbf{Audience} & \textbf{Type} \\
\midrule
VotaSJ PWA \& audit portal & Citizens, journalists & Own digital \\
The Graph subgraph + REST API & TCE-SC, MP, academia & Own M2M \\
gov.br OIDC integration page & Citizens (onboarding) & Partner digital \\
Assisted kiosks at CRAS/UPAs & Unbanked, elderly & Own physical \\
Neighborhood associations & Grassroots citizens & Partner grassroots \\
Regional press (NDTV/SC, ND) & All citizens & Partner media \\
Conferences (SBSeg, SBRC, LADC) & Auditors, researchers & Own academic \\
Open-source repository & Engineers, municipal IT & Own technical \\
\bottomrule
\end{tabular}
\end{center}

\subsection{Content artifacts}

The white-paper series itself is the spine of the strategy. Three
artifacts, deliberately layered: (1) this \textbf{long technical paper}
for engineers, CTOs, TCE-SC technical staff, and academic reviewers;
(2) a \textbf{short business/user paper} (companion deliverable, leaflet
style) for elected officials and citizens; and (3) a \textbf{pitch deck}
and \textbf{demo video} for funding-oriented audiences. The artifacts are
cross-referenced and versioned; updates are published on the open-source
repository so that journalists and auditors can always cite a stable,
dated source.

\subsection{Pilot-driven activation}

The first municipal cycle is the marketing event of the year. Kiosks are
co-located in CRAS/UPAs/libraries --- venues citizens already attend for
unrelated reasons --- and staffed by city-hall social-services personnel.
The cycle produces three by-products of equal commercial value:
(a) measured turnout (the only metric a competitor cannot fake);
(b) photographic evidence of execution milestones (the substance of the
``accountability loop'' value proposition); and (c) press-ready stories of
named citizens whose votes turned into named municipal works.

\section{Future Roadmap}
\label{sec:roadmap}

\noindent\begin{tabular}{l l p{0.45\textwidth}}
\toprule
\textbf{Phase} & \textbf{Target} & \textbf{Scope} \\
\midrule
Delivered & Contract suite + spec & \texttt{VoterRegistry}, \texttt{ParticipatoryBudget} and \texttt{ExecutionTracker} implemented, 40-test Hardhat suite, gas envelope characterized. \\
2026-Q3 & Field pilot & Single cycle in one administrative region of S\~ao Jos\'e, 5--10k participants. \\
2026-Q4 & Full S\~ao Jos\'e cycle & Citywide, \(\approx\)50k voters. \\
2027 & Multi-municipality & Onboarding Palho\c{c}a, Bigua\c{c}u, Florian\'opolis. \\
2027--2028 & zk-Semaphore default & Full vote-secrecy from City Hall, not just from the public~\cite{semaphoredocs}. \\
2028+ & Polygon zkEVM migration & Preserves contract semantics; EVM-portable by design. \\
\bottomrule
\end{tabular}

\section{Risks, Limitations, and Ethical Considerations}
\label{sec:risks}

What VotaSJ does \textbf{not} solve:

\begin{itemize}
  \item It does not produce deliberation. A digital ballot is not an assembly.
  \item It does not solve the digital divide. Roughly 8\,\% of S\~ao Jos\'e
        adults do not own a smartphone; assisted kiosks mitigate but do not
        eliminate the exclusion.
  \item It is not a coup-proof system. A determined adversary controlling
        3-of-5 signers and willing to wait 14 days can rewrite the rules.
  \item Polygon is a dependency. If Polygon ceases to be credible, migration
        is non-trivial (though tractable, given EVM portability).
  \item Legal feasibility of binding digital voting under the S\~ao Jos\'e
        Organic Law is not formally established. The first field pilot is
        consultative, not binding, pending an OAB-SC and municipal legal
        opinion.
\end{itemize}

\noindent A few constraints we treat as non-negotiable. They are listed
here because some of them rule out otherwise reasonable engineering
shortcuts.

\begin{itemize}
  \item \textbf{Citizens pay nothing --- money, gas, or cognitive overhead.}
        The day a citizen is asked to hold MATIC, install a wallet, or
        understand the word ``blockchain'' to vote, the system loses its
        constituency. The ERC-4337~\cite{eip4337} paymaster, the in-app
        wallet abstraction, and the gov.br SSO are there to keep that
        promise, not to be impressive.
  \item \textbf{No tradeable representation of voting power.} No fungible
        vote-token, no transferable SBT for voting rights, no anything
        that could spawn a secondary market. Participation badges are
        non-transferable (ERC-5192-style
        SBT~\cite{eip5192,ohlhaver2022desoc}), per-cycle, wallet-bound.
  \item \textbf{Pseudonymity, not anonymity --- and we say so.} Today's
        votes link wallet $\leftrightarrow$ \texttt{credentialHash},
        opaque to the public but visible to the City Hall via the
        registry. The second annual cycle moves to Semaphore-based
        anonymity. The paper does not claim a property the system does
        not yet enforce.
  \item \textbf{Source available, audit available.} Contracts, indexer
        schema, and audit reports are public. No ``secret sauce''.
  \item \textbf{No surveillance side-channel.} Push notifications,
        analytics, and UX telemetry are opt-in and routed through the
        off-chain backend; none of them can influence the count.
\end{itemize}

\section{Conclusion}
\label{sec:conclusion}

Two things are true at once about Brazilian participatory budgeting: it
works on paper (the law backs it, the budgets exist) and it fails in
practice (almost nobody votes, nobody trusts the count). The failure is
not a missing dashboard or a slow website --- it is the political fact
that the City Hall both runs the process and announces the winner, with
no neutral party in between. A blockchain does not fix democracy, but it
does fix that particular asymmetry: a public L2 with an external
validator set~\cite{polygondocs} produces a count that nobody in S\~ao
Jos\'e has to take on faith.

VotaSJ uses that property as narrowly as possible. Eligibility, votes,
tally, and execution milestones live on-chain; everything else stays
where it belongs --- \texttt{gov.br}~\cite{govbrmanual} for identity,
IPFS for proposal content, a normal web backend for UX. The 3-of-5
multisig with two external signers and the 14-day
timelock~\cite{openzeppelin5} make it impossible for any single party,
including the City Hall itself, to change the rules during an active
vote. The per-cycle cost lands at roughly one-third of the existing
process. The citizen never has to know what a blockchain is.

What remains open is institutional, not technical: a legal opinion on
the bindingness of digital voting under the S\~ao Jos\'e Organic Law,
an ANPD position on hash-anchored eligibility, and the operational
readiness of the City Hall as the anchor customer. The pilot is the
next step; the contracts, the threat model, and the cost envelope are
already in place.

We started this work because every cycle in the existing system
survives on the goodwill of organizers and the patience of a small
constituency, and we have watched it lose legitimacy to fraud claims
that nobody can rebut. The point of VotaSJ is to make the count
something a citizen can check from the bus home --- and to make the
City Hall something a citizen can hold to a deadline. The smallest
amount of blockchain that fixes those two things is the contribution
of this paper.

\bibliographystyle{sbc}
\bibliography{refs}

\end{document}
